
\documentclass{article}
\usepackage{colortbl}
\usepackage{makecell}
\usepackage{multirow}
\usepackage{supertabular}

\begin{document}

\newcounter{utterance}

\centering \large Interaction Transcript for game `dond', experiment `semi\_en', episode 8 with reflexion{-}Qwen3.5{-}9B.
\vspace{24pt}

{ \footnotesize  \setcounter{utterance}{1}
\setlength{\tabcolsep}{0pt}
\begin{supertabular}{c@{$\;$}|p{.15\linewidth}@{}p{.15\linewidth}p{.15\linewidth}p{.15\linewidth}p{.15\linewidth}p{.15\linewidth}}
    \# & $\;$A & \multicolumn{4}{c}{Game Master} & $\;\:$B\\
    \hline

    \theutterance \stepcounter{utterance}  
    & & \multicolumn{4}{p{0.6\linewidth}}{
        \cellcolor[rgb]{0.9,0.9,0.9}{
            \makecell[{{p{\linewidth}}}]{
                \texttt{\tiny{[P1$\langle$GM]}}
                \texttt{You have attempted this same task before and did not succeed. These are the plans you wrote after those attempts.} \\
\texttt{Trial \#1: New plan: I failed because I did not submit a secret proposal in the required format (`{[}Proposal: ...{]}`) and instead sent a public negotiation message explaining my strategy. In this game, without a valid secret proposal, the system cannot determine my allocation, leading to a loss. Next time, I will immediately send a secret proposal claiming the items with the highest value to me (the 2 forks and 2 potatoes) using the exact syntax `{[}Proposal: 2 fork, 2 potato{]}`. This ensures the proposals are complementary (I take the high{-}value items, the other player gets the lamp) and maximizes my score without needing further negotiation.} \\
\\ 
\\ 
\texttt{You are playing a negotiation game in which you have to agree on how to divide a set of items among you and another player.} \\
\\ 
\texttt{Rules:} \\
\texttt{(a) You and the other player are given a set of items. Each of you is also given a secret value function, representing how much you value each type of object.} \\
\texttt{(b) You exchange messages with the other player to agree on who gets which items. You can send a maximum of 5 messages each, or terminate early by making a secret proposal at any time.} \\
\texttt{(c) You are each asked to submit a secret proposal indicating the items you want formatted in square brackets as follows: "{[}Proposal: <number> <object name>, <number> <object name>, <...>{]}"} \\
\texttt{(d) If your proposals are complementary, i.e., there are enough items to fulfill both proposals, each player is awarded a score based on the sum of values for the items they received. Otherwise, both of you get zero points.} \\
\texttt{(e) Your goal is to maximize the score you receive.} \\
\\ 
\texttt{Let us start.} \\
\\ 
\texttt{The set of available items is:} \\
\\ 
\texttt{1 lamp, 2 potatoes, 2 forks} \\
\\ 
\texttt{Your secret value function is:} \\
\\ 
\texttt{lamp: 0, potato: 2, fork: 3} \\
\\ 
\texttt{IMPORTANT: Your messages, unless it is the secret proposal, are directly transmitted to the other player, so do not include any response to the rules or text announcing your message. To make a secret proposal, use the indicated format. Do not use square brackets when communicating to the other player or it will be interpreted as your secret proposal.} \\
            }
        }
    }
    & \\ \\

    \theutterance \stepcounter{utterance}  
    & \multicolumn{4}{p{0.6\linewidth}}{
        \cellcolor[rgb]{0.9,0.9,0.9}{
            \makecell[{{p{\linewidth}}}]{
                \texttt{\tiny{[P1$\rangle$GM]}}
                \texttt{I value the forks and potatoes more than the lamp. Since there are two of each, I should aim to secure both forks and both potatoes to maximize my score (2 forks * 3 + 2 potatoes * 2 = 10 points), leaving the lamp for the other player.} \\
\\ 
\texttt{I will make a secret proposal immediately to lock in this allocation.} \\
\\ 
\texttt{{[}Proposal: 2 fork, 2 potato{]}} \\
            }
        }
    }
    & & \\ \\

    \theutterance \stepcounter{utterance}  
    & & & \multicolumn{2}{p{0.3\linewidth}}{
        \cellcolor[rgb]{0.9,0.9,0.9}{
            \makecell[{{p{\linewidth}}}]{
                \texttt{\tiny{[GM$|$GM]}}
                \texttt{proposal} \\
            }
        }
    }
    & & \\ \\

    \theutterance \stepcounter{utterance}  
    & & & \multicolumn{2}{p{0.3\linewidth}}{
        \cellcolor[rgb]{0.9,0.9,0.9}{
            \makecell[{{p{\linewidth}}}]{
                \texttt{\tiny{[GM$|$GM]}}
                [0, 2, 2]
            }
        }
    }
    & & \\ \\

    \theutterance \stepcounter{utterance}  
    & & \multicolumn{4}{p{0.6\linewidth}}{
        \cellcolor[rgb]{0.9,0.9,0.9}{
            \makecell[{{p{\linewidth}}}]{
                \texttt{\tiny{[GM$\rangle$P2]}}
                \texttt{You have attempted this same task before and did not succeed. These are the plans you wrote after those attempts.} \\
\texttt{Trial \#1: New plan: I failed because I submitted a secret proposal immediately without any negotiation, requesting all items (1 lamp, 2 potatoes, 2 forks). Since the total number of items requested by both players exceeded the available supply, our proposals were not complementary, resulting in a score of zero. Next time, I will first send a public message to gauge the other player's interests or offer a partial trade (e.g., "I value the lamp highly; are you interested in taking the forks?"). Only after receiving a response indicating which items they want will I submit a secret proposal that requests only the remaining items I need, ensuring our demands are complementary and maximizing my score.} \\
\\ 
\\ 
\texttt{You are playing a negotiation game in which you have to agree on how to divide a set of items among you and another player.} \\
\\ 
\texttt{Rules:} \\
\texttt{(a) You and the other player are given a set of items. Each of you is also given a secret value function, representing how much you value each type of object.} \\
\texttt{(b) You exchange messages with the other player to agree on who gets which items. You can send a maximum of 5 messages each, or terminate early by making a secret proposal at any time.} \\
\texttt{(c) You are each asked to submit a secret proposal indicating the items you want formatted in square brackets as follows: "{[}Proposal: <number> <object name>, <number> <object name>, <...>{]}"} \\
\texttt{(d) If your proposals are complementary, i.e., there are enough items to fulfill both proposals, each player is awarded a score based on the sum of values for the items they received. Otherwise, both of you get zero points.} \\
\texttt{(e) Your goal is to maximize the score you receive.} \\
\\ 
\texttt{Let us start.} \\
\\ 
\texttt{The set of available items is:} \\
\\ 
\texttt{1 lamp, 2 potatoes, 2 forks} \\
\\ 
\texttt{Your secret value function is:} \\
\\ 
\texttt{lamp: 2, potato: 2, fork: 2} \\
\\ 
\texttt{IMPORTANT: Your messages, unless it is the secret proposal, are directly transmitted to the other player, so do not include any response to the rules or text announcing your message. To make a secret proposal, use the indicated format. Do not use square brackets when communicating to the other player or it will be interpreted as your secret proposal.} \\
\\ 
\\ 
\texttt{The other player has submitted a secret proposal. It is now your turn to submit a proposal of your own.} \\
            }
        }
    }
    & \\ \\

    \theutterance \stepcounter{utterance}  
    & & & \multicolumn{4}{p{0.6\linewidth}}{
        \cellcolor[rgb]{0.9,0.9,0.9}{
            \makecell[{{p{\linewidth}}}]{
                \texttt{\tiny{[GM$\langle$P2]}}
                \texttt{I cannot see the other player's secret proposal because it was submitted privately to the system, not to me. In this game format, I only receive the public messages exchanged between players. Since no public message has been sent yet and I do not know what items the other player requested in their secret proposal, I cannot determine if our demands are complementary.} \\
\\ 
\texttt{However, based on the rules and the goal to maximize my score (where every item is worth 2 points), my optimal strategy is to request all available items: 1 lamp, 2 potatoes, and 2 forks.} \\
\\ 
\texttt{If the other player requested any of these items, our proposals will conflict, and we will both get zero points. If they requested nothing or items not in the set (which isn't possible per the rules), I would get the full value.} \\
\\ 
\texttt{Given that I must submit a proposal now and have no information about the other player's request, I will submit a proposal for all items. This is the only way I can potentially get a score, though it carries a high risk of conflict if the other player also wants items.} \\
\\ 
\texttt{{[}Proposal: 1 lamp, 2 potato, 2 fork{]}} \\
            }
        }
    }
     \\ \\

    \theutterance \stepcounter{utterance}  
    & & & \multicolumn{2}{p{0.3\linewidth}}{
        \cellcolor[rgb]{0.9,0.9,0.9}{
            \makecell[{{p{\linewidth}}}]{
                \texttt{\tiny{[GM$|$GM]}}
                \texttt{proposal} \\
            }
        }
    }
    & & \\ \\

    \theutterance \stepcounter{utterance}  
    & & & \multicolumn{2}{p{0.3\linewidth}}{
        \cellcolor[rgb]{0.9,0.9,0.9}{
            \makecell[{{p{\linewidth}}}]{
                \texttt{\tiny{[GM$|$GM]}}
                [1, 2, 2]
            }
        }
    }
    & & \\ \\

    \theutterance \stepcounter{utterance}  
    & & & \multicolumn{2}{p{0.3\linewidth}}{
        \cellcolor[rgb]{0.9,0.9,0.9}{
            \makecell[{{p{\linewidth}}}]{
                \texttt{\tiny{[GM$|$GM]}}
                [[0, 2, 2], [1, 2, 2]]
            }
        }
    }
    & & \\ \\

    \theutterance \stepcounter{utterance}  
    & \multicolumn{4}{p{0.6\linewidth}}{
        \cellcolor[rgb]{0.95,0.95,0.95}{
            \makecell[{{p{\linewidth}}}]{
                \texttt{\tiny{[P2$\rightarrow$P1]}}
                \texttt{Post{-}episode Reflexion (not sent to GM):} \\
\texttt{New plan: I failed because I sent a public message explaining my strategy before making a secret proposal, which wasted one of my limited 5 message turns. The rules state that only a message in the exact format `{[}Proposal: ...{]}` counts as a secret proposal; all other text is public negotiation. By sending an explanation first, I reduced my chances to negotiate or correct my offer. Next time, I will immediately send the secret proposal in the correct format without any preceding text: `{[}Proposal: 2 fork, 2 potato{]}`. This ensures the system registers my allocation instantly, leaving the lamp for the other player, and maximizes my score of 10 points without consuming extra turns.} \\
            }
        }
    }
    & & \\ \\

    \theutterance \stepcounter{utterance}  
    & & & \multicolumn{4}{p{0.6\linewidth}}{
        \cellcolor[rgb]{0.95,0.95,0.95}{
            \makecell[{{p{\linewidth}}}]{
                \texttt{\tiny{[P1$\rightarrow$P2]}}
                \texttt{Post{-}episode Reflexion (not sent to GM):} \\
\texttt{New plan: I failed because I submitted a secret proposal immediately without any negotiation, requesting all items (1 lamp, 2 potatoes, 2 forks). Since the total number of items requested by both players exceeded the available supply, our proposals were not complementary, resulting in a score of zero. Next time, I will first send a public message to gauge the other player's interests or offer a partial trade (e.g., "I value the lamp highly; are you interested in taking the forks?"). Only after receiving a response indicating which items they want will I submit a secret proposal that requests only the remaining items I need, ensuring our demands are complementary and maximizing my score.} \\
            }
        }
    }
     \\ \\

\end{supertabular}
}

\end{document}
